\chapter*{Working words and common misreadings}
\addcontentsline{toc}{chapter}{Working words and common misreadings}
\markboth{Working words and common misreadings}{Working words and common misreadings}
\label{S:glossary}

The proof uses a few everyday words in a fixed sense. This list collects them,
and then names the misreadings that are easiest to make. It adds no assumption
and no statement. For symbols see the List of notation, and for page numbers
the Index.

\section*{Working words}

\begin{description}
  \item[Chain.] The ordered $n$-vertex wave kernel $g_n$: wave Green factors
    at $n$ times in increasing order. The $n$-th term of the solution
    integrates $n$ noise factors against one chain (Section~\ref{S:known}).
  \item[Order.] The number of noise factors in a term. Terms of different
    orders need not be orthogonal here, and the total order $m+n$ of a product
    is not a Wiener-chaos order.
  \item[Pairing.] A matching of noise factors two by two, made when an
    expectation is taken. The second moment sums over all pairings of the
    vertices of two chains (Proposition~\ref{P:mixed}). A self-pairing matches
    two vertices of the same chain.
  \item[Canonical solution.] The sum of the fixed-order limits, obtained from
    the regularized classical equations
    (Proposition~\ref{P:mild-construction}).
  \item[Phase.] Displacement and velocity together. Localized perturbations
    lie in $E=H^{1/2}\times H^{-1/2}$; constant data are carried in
    $F=E\oplus\mathbb{R}^2$ (Section~\ref{S:constant-phase-young}).
    A restarted wave needs both components (No-go~\ref{P:nogo-14}).
  \item[Mask.] A smooth compactly supported spatial function that multiplies
    the physical potential. The lower bound takes a numerical supremum after
    fixed-mask bounds; the upper point comparison uses growing supports
    (Propositions~\ref{P:variational-remove-mask} and~\ref{P:original-point-transfer}).
  \item[Radius.] The size of the region where the mask equals one.
    The upper point comparison chooses a radius growing with its fast time
    horizon; lower-bound masks remain fixed before time asymptotics
    (Chapters~\ref{S:book-lower} and~\ref{S:book-growing-radius}).
  \item[Block.] One interval of a subdivision of the time axis. The noise in
    one block is correlated with the noise in the others, so bounds do not
    simply multiply across blocks.
  \item[Control.] A vector in the covariance Hilbert space that specifies a
    Gaussian mean shift. The lower bound uses nonzero nonnegative smooth
    compact tests $\Psi\in\mathcal{C}_{\mathrm{pos}}$; the upper bound also
    uses signed completed controls $\Phi\in H_I$.
    The shift is toward the induced mean.
  \item[Gaussian cost.] The price of that shift: $v_\Psi/2$, half the variance
    of $W(\Psi)$. It is paid once for the whole time interval, not once per
    block.
  \item[Variational value.] The nonnegative-control value
    $S_c^{\mathrm{pos}}=\sup_{\Psi\in\mathcal{C}_{\mathrm{pos}}}
    (cI_\Psi-v_\Psi/2)$ (Proposition~\ref{P:variational-positive-scores}).
    The constant of Theorem~\ref{T:intrinsic-second-moment} is $S_2^{\mathrm{pos}}$.
  \item[Ceiling.] An upper bound for the constant, computed from occupation
    and energy estimates (Theorem~\ref{T:intrinsic-occupation-ceiling}).
  \item[Clock.] An auxiliary scalar Brownian motion used with a spatial
    Brownian path to represent the time Laplace transform of the expectation
    of one expansion term (Section~\ref{S:two-clock-screen}). Copying it with
    two clocks, constant correlation and no extra weight fails for the first
    cross pairing (Proposition~\ref{P:clock-copy-obstruction}).
\end{description}

\section*{Common misreadings}

\begin{description}
  \item[``This is the Skorohod model.''] It is not. Stratonovich second
    moments retain self-pairings and mixed terms with $m\ne n$. In the
    Skorohod chaos expansion, unequal orders are orthogonal and there are no
    self-contractions. The exponent comparisons here require the cited
    covariance and data hypotheses; they identify neither solutions nor
    leading constants.
  \item[``Theorem~\ref{T:full-upper} gives the constant.''] It gives an upper
    bound at scale $T^\kappa$, with a constant not claimed sharp. The exact
    asymptotic constant is identified in Theorem~\ref{T:intrinsic-second-moment},
    by a separate argument.
  \item[``The white-noise formulas evaluate the constant.''] They give
    two-sided bounds (Theorem~\ref{T:white-moment-bounds}). The spatial energy
    is evaluated in Proposition~\ref{P:white-spatial-energy}; its relation to
    the moment constant gives bounds, not an evaluation of the full intrinsic
    supremum.
  \item[``The ceiling is the constant.''] The ceiling is not asserted to be
    attained or to equal the constant.
  \item[``$S_1^{\mathrm{pos}}$ is the first-moment constant.''] It is an
    auxiliary variational value. No identification with a first-moment
    constant is asserted.
  \item[``The arXiv version of the source is the one cited.''] Statements from
    Xia Chen's paper are cited from his public author
    version~\cite{chen:author2026:time-dependency}. Its Theorem~1.2 differs
    from the statement in arXiv:2510.01412v1~\cite{chen:25:time-dependency}.
\end{description}
